% Generated draft: requires mathematical and source review.
\documentclass[11pt]{article}
\usepackage[margin=25mm]{geometry}
\usepackage{fontspec}
\setmainfont{Latin Modern Roman}
\usepackage{amsmath,amssymb,mathtools,unicode-math}
\setmathfont{Latin Modern Math}
\usepackage{enumitem,needspace,xurl,hyperref}
\setlength{\parindent}{0pt}
\setlength{\parskip}{5pt}
\setlength{\emergencystretch}{3em}
\newcommand{\N}{\mathbb{N}}
\newcommand{\Z}{\mathbb{Z}}
\newcommand{\Q}{\mathbb{Q}}
\newcommand{\R}{\mathbb{R}}
\newcommand{\C}{\mathbb{C}}
\newcommand{\F}{\mathbb{F}}
\newcommand{\A}{\mathbb{A}}
\newcommand{\PP}{\mathbb{P}}
\newcommand{\E}{\mathbb{E}}
\newcommand{\eps}{\varepsilon}
\newcommand{\dd}{\,\mathrm{d}}
\begin{document}
\section*{UnsolvedMath ID: 701658}
\textbf{Source problem code:} catalogue entry
\subsection{Definitions and mathematical statement}
A variety of groups is a class \(\mathcal V\) consisting of all groups satisfying some fixed set of group identities. For a positive integer \(n\), the relatively free group \(F_{\mathcal V}(n)\) of rank \(n\) is a group in \(\mathcal V\) with distinguished generators \(x_1,\ldots,x_n\) such that, for every \(G\in\mathcal V\) and every \((g_1,\ldots,g_n)\in G^n\), there is a unique homomorphism \(\varphi:F_{\mathcal V}(n)\to G\) satisfying \(\varphi(x_i)=g_i\) for all \(1\leq i\leq n\).

Does there exist a variety of groups \(\mathcal V\) such that
\[
|F_{\mathcal V}(2)|<\aleph_0
\qquad\text{and}\qquad
|F_{\mathcal V}(3)|\geq\aleph_0?
\]
The supplied record labels this question open; that label is not an independently verified mathematical or literature assessment. The supplied attribution names P. Zusmanovich as proposer and also places G. M. Bergman and A. Magidin before the question, without explaining their respective roles.
\subsection{Short English statement}
Can a variety of groups have a finite relatively free group on two generators but an infinite relatively free group on three generators?
\subsection{Sources}
\begin{itemize}
\item Local UnsolvedMath registry, permanent ID \(701658\), source problem code \texttt{catalogue entry}, title \emph{Kourovka Notebook Problem 19.110}. The supplied record carries the status \texttt{open}.
\item Supplied source attribution: \emph{The Kourovka Notebook (Unsolved Problems in Group Theory)}, 19th issue (2018), problem 19.110; also identified in the record with the 21st issue edition (2026), \texttt{arXiv:1401.0300v48}. The collection description names E. I. Khukhro and V. D. Mazurov as editors. These bibliographic details are supplied source data and have not been independently verified here.
\item The background and proposer field credit P. Zusmanovich, with proposed year 2018. The statement additionally begins with the parenthetical attribution ``G. M. Bergman and A. Magidin''; the record does not explain the relationship between these credits.
\end{itemize}
\end{document}
